\section{Tres familias de modelos sobre el mismo esqueleto}

La sección anterior implementó un GPT decoder-only. En esta sección solo se cambian la connectivity y el objective, y así se obtienen las familias encoder-only y encoder--decoder. Esta perspectiva unificada es más confiable que memorizar nombres de modelos: primero se pregunta a quién puede ver cada posición, después de dónde vienen las keys/values y, por último, qué salida exige el objetivo de entrenamiento.

\subsection{Encoder: se elimina la causal mask}

La self-attention del encoder permite que cada token lea al mismo tiempo el contexto de la izquierda y el de la derecha, por lo que sirve para generar una contextual representation, para clasificación o para denoising. En el código basta con eliminar la causal mask para que todas las posiciones de la attention matrix participen en el softmax; en cambio, el objetivo ya no puede ser la misma next-token loss left-to-right, porque los token futuros filtrarían la respuesta.

\lecturefigure{slide-44-encoder-blocks.jpg}{Encoder blocks: se elimina el causal masking para que todos los token se comuniquen en ambas direcciones.}{Video oficial de Stanford Online, 00:43:18--00:43:53.}

\begin{knowledgebox}{Lectura de la figura: la misma attention, con restricciones distintas}
La connectivity triangular inferior de GPT impone la generación causal; la self-attention totalmente conectada de los encoder tipo BERT permite que cada token reúna la información de toda la oración. BERT suele usar masked/denoising objectives: primero se corrompe la entrada y luego se pide al modelo que la recupere o que extraiga una representación. La diferencia de arquitectura no se reduce a «si hay o no mask»: el objective también debe corresponder a la información visible.
\end{knowledgebox}

\subsection{Cross-attention: el decoder consulta la encoder memory}

Un modelo encoder--decoder primero codifica la secuencia fuente como $H_{\text{enc}}$; el decoder state $X_{\text{dec}}$ produce las queries, mientras que las keys/values provienen del encoder:

\[
Q=X_{\text{dec}}W_Q,\qquad
K=H_{\text{enc}}W_K,\qquad
V=H_{\text{enc}}W_V.
\]

Así, cada token de destino puede ver el prefijo ya generado mediante la causal self-attention y, además, leer la entrada fuente completa mediante la cross-attention. La traducción automática, la generación de resúmenes y las text-to-text tasks pueden aprovechar este canal de condicionamiento explícito.

\lecturefigure{slide-45-cross-attention.jpg}{Cross-attention: las decoder queries leen la información fuente a partir de las encoder keys/values.}{Video oficial de Stanford Online, 00:43:53--00:44:47.}

\begin{knowledgebox}{Lectura de la figura: la dirección de las flechas indica el origen de la información}
En la decoder self-attention, Q/K/V provienen de los decoder states; en la cross-attention, la query sigue saliendo del decoder node actual, pero la key/value proviene de la capa superior del encoder. Ambas comparten el dot product, el softmax y la value aggregation. Lo «cross» no es otro algoritmo misterioso, sino la misma operación de direccionamiento aplicada entre dos conjuntos de states.
\end{knowledgebox}

\subsection{Comparación estructurada de GPT, BERT y T5}

Cuando hay muchos nombres de modelos, conviene reducirlos a su connectivity y su objective. La tabla siguiente describe solo los prototipos clásicos; los modelos reales pueden incorporar span corruption, prefix LM, mixture-of-experts o multimodal adapters, pero al evaluar un modelo nuevo se pueden seguir aplicando las mismas tres preguntas.

\lecturefigure{slide-46-transformer-model-families.jpg}{Las tres grandes familias: GPT decoder-only, BERT encoder-only y T5 encoder--decoder.}{Video oficial de Stanford Online, 00:44:47--00:45:31.}

\begin{table}[H]
\centering
\small
\begin{tabularx}{\textwidth}{L{2.8cm}L{3.5cm}L{3.2cm}X}
\toprule
Familia & Connectivity & Objetivo típico & Problemas adecuados \\
\midrule
Decoder-only / GPT & Autoatención causal de izquierda a derecha & Predecir el siguiente token & Generación abierta, completar a partir de un prompt \\
Encoder-only / BERT & Autoatención bidireccional & Recuperación de lo enmascarado, eliminación de ruido o tareas supervisadas & Extracción de representaciones, clasificación, recuperación de información, etiquetado de secuencias \\
Encoder--decoder / T5 & Codificador totalmente conectado; decodificador causal; atención cruzada & Generación condicionada de secuencias & Traducción, resúmenes, transformaciones estructuradas \\
\bottomrule
\end{tabularx}
\caption{Términos clave: una familia de modelos se define por la información visible y por el objetivo de entrenamiento en conjunto.}
\end{table}

\teachervoice{En la sesión de Q\&A, Karpathy dijo que no le gusta mucho el proceso autorregresivo en el que cada token muestreado queda comprometido de forma permanente, y que prefiere algo parecido a diffusion: «primero un borrador y luego una revisión». Se trata de una intuición personal de investigación de 2023, no de una dirección de evolución inevitable del Transformer demostrada en esta clase.}

\subsection{Resumen de la sección}

Decoder-only usa causal self-attention para la next-token generation; encoder-only permite la comunicación bidireccional y adopta objetivos de aprendizaje de representaciones que no filtran las etiquetas; encoder--decoder hace que un causal decoder consulte la memory de la fuente mediante cross-attention. Las tres comparten las matemáticas de la attention y solo se reconfiguran en la connectivity, la state source y el objective.
